%% Rechnungsvorlage „helsinki“ (ADR-0011): ein Blickfang, der Titelbalken, sonst ruhiger Textfluss; Linien nur,
%% wo sie Zahlen gliedern. Jede einzelne Angabe darf leer sein, entschieden wird mit \ifblank und nichts sonst
%% (ADR-0017).
\documentclass[11pt,a4paper]{article}

%% Als Erstes, denn die Grundlage lädt pdfx; was sie lädt, lädt diese Vorlage nicht noch einmal (ADR-0022).
\usepackage{grundlage}

\usepackage{geometry}
%% colortbl für \arrayrulecolor; xcolor lädt die Grundlage.
\usepackage{colortbl}
\usepackage{booktabs}  %% \midrule bringt die vertikalen Abstände selbst mit
%% xltabular lädt tabularx und longtable selbst; tabularx bleibt für Summen und Absetzung (ADR-0020).
\usepackage{xltabular}
\usepackage{fancyhdr}
\usepackage{array}
\usepackage{etoolbox}  %% \ifblank für die Angaben, die leer sein dürfen
\usepackage{tikz}      %% das Logo

%% Gruppiert ab vier Stellen; Trennzeichen und Stellung des Währungszeichens setzt die Sprachdatei (ADR-0021).
\sisetup{group-minimum-digits = 4}

%% Das Zeichen dazu setzt die Sprachdatei in \betrag (ADR-0026).
\newcommand{\waehrung}{<<WAEHRUNG>>}

%% Nach etoolbox und siunitx, denn die Sprachdatei braucht beide.
\input{sprache}

%% Dieselbe Schrift wie im Lebenslauf – dadurch passen die Dokumente zusammen.
\setmainfont{texgyreheros}[
  Extension      = .otf,
  UprightFont    = *-regular,
  BoldFont       = *-bold,
  ItalicFont     = *-italic,
  BoldItalicFont = *-bolditalic,
]

%% Die Vorgabe der Akzentfarbe steht hier und nicht in Kotlin (ADR-0024): das Mittel aus dem Rot der Webseite
%% und dem Ton der Rechnung vor dem Generator, beide Ausgangsfarben stehen daneben.
\ifblank{<<AKZENTFARBE>>}{\definecolor{akzent}{HTML}{8D313A}}{\definecolor{akzent}{HTML}{<<AKZENTFARBE>>}}
\definecolor{akzentweb}{HTML}{8B1A2B}
\definecolor{akzentalt}{HTML}{8E4848}
\definecolor{haupttext}{HTML}{1C1C1E}
\definecolor{gedaempft}{HTML}{5C5C62}
\definecolor{linie}{HTML}{C9C9CF}

%% Außerhalb der Tabelle: Darin eröffnete \arrayrulecolor eine Zelle, und das \noalign von \midrule landete falsch.
\arrayrulecolor{linie}

%% includefoot misst den unteren Rand bis zur letzten Grundlinie des Fußes. Der Kopf der Folgeseiten steht im
%% oberen Rand, damit die erste Seite ihn nicht mitbezahlt.
\geometry{a4paper, left=25mm, right=20mm, top=18mm, bottom=16mm,
  includefoot, footskip=12mm, headheight=14pt, headsep=6mm}

\color{haupttext}
\setlength{\parindent}{0pt}
%% Folgeseiten tragen einen kurzen Kopf, damit eine lose Seite zu ihrer Rechnung findet.
\fancypagestyle{erste}{\fancyhf{}\renewcommand{\headrulewidth}{0pt}\fancyfoot[C]{\fuss}\fancyfoot[R]{\fortsetzung}}
\fancypagestyle{folge}{\fancyhf{}\renewcommand{\headrulewidth}{0pt}%
  \fancyhead[L]{\small\color{gedaempft}<<DOKUMENTART>> <<RECHNUNGSNUMMER>> · <<RECHNUNGSDATUM>>}%
  \fancyhead[R]{\small\color{gedaempft}\wortseitevon{\thepage}{\pageref*{LastPage}}}%
  \fancyfoot[C]{\fuss}\fancyfoot[R]{\fortsetzung}}
\pagestyle{folge}

%% Im Fuß, damit er an jeder Seite hängt, ohne Raum einzunehmen: 3 mm unter dem Ende des Textbereichs.
\newcommand{\fortsetzung}{\wennweitereseite{%
  \raisebox{\dimexpr\footskip-3mm\relax}[0pt][0pt]{\footnotesize\color{gedaempft}\wortfortsetzung}}}

%% Flattersatz: In den schmalen Spalten risse Blocksatz Löcher zwischen die Wörter.
\raggedright

%% Keine Silbentrennung: Getrennte Eigennamen und Firmennamen sehen falsch aus.
\hyphenpenalty=10000
\exhyphenpenalty=10000
\tolerance=2000

%% Achtung: minipage setzt über \@parboxrestore die Absatzparameter zurück, der Flattersatz gilt dort nicht; ein
%% \AtBeginEnvironment-Hook kommt zu früh. Deshalb steht \raggedright in jeder Minipage mit Fließtext.

\newlength{\balkenhoehe}

%% Das Logo definiert \logozeichnung und optional \logoskalierung (ADR-0012). TikZ und keine Bilddatei:
%% luatex.def gibt jedem per \includegraphics eingebundenen PDF eine Transparenzgruppe mit, und die verbietet
%% PDF/A-1 (#197).
\newlength{\logohoehe}\setlength{\logohoehe}{12mm}
\newlength{\versalhoehe}
\ifblank{<<LOGO>>}{}{\input{<<LOGO>>}}
\providecommand{\logoskalierung}{1}

%% Abstände als benannte Längen, damit das Satzbild an einer Stelle steht (#51). 2.5mm zwischen Positionen:
%% Bei knappem Abstand verschmölze eine Position mit Zusatzzeile optisch mit der nächsten.
\newlength{\positionsabstand}\setlength{\positionsabstand}{2.5mm}
\newlength{\summenabstand}\setlength{\summenabstand}{1mm}
\newlength{\blockabstand}\setlength{\blockabstand}{3mm}
\newlength{\ustidabstand}\setlength{\ustidabstand}{2mm}

%% Der Renderer ruft nur das Makro auf, nie die Länge: Ein vertippter Längenname fiele erst LuaLaTeX auf (#76).
\newcommand{\positionstrenner}{\\[\positionsabstand]}
\newcommand{\summentrenner}{\\[\summenabstand]}
\newcommand{\blocktrenner}{\\[\blockabstand]}
\newcommand{\ustidtrenner}{\\[\ustidabstand]}

%% Aufgerufen als \zeilenwechsel{}: Das leere Paar schützt die nächste Zeile, denn beginnt sie mit „[“, läse \\
%% sie als Abstandsangabe. Ein rohes \\ steht deshalb nur vor festem Text, nie vor einem Platzhalter (#147).
\newcommand{\zeilenwechsel}{\\}

%% In Versalien nach DIN 5008 (#374).
\newcommand{\landzeile}[1]{\MakeUppercase{#1}}

%% Das {} hinter \\ wie beim Aufruf von \zeilenwechsel.
\newcommand{\kopfperson}[1]{\\{}#1}

%% Ist der Wert leer, entfällt die Zeile samt Beschriftung. \par statt \\: Ein \\ hinter der letzten Zeile ließe
%% eine leere Zeile zurück.
\newcommand{\kopfangabe}[2]{\ifblank{#2}{}{#1 #2\par}}

%% Im schmalen Block bräche „Juni bis Oktober 2026“ sonst hinter dem „bis“.
\newcommand{\zusammen}[1]{\mbox{#1}}
%% \betrag und \prozent stehen in der Sprachdatei.
\newcommand{\menge}[1]{\num{#1}}

%% \rubrik statt \label, das in LaTeX vergeben ist.
\newcommand{\rubrik}[1]{{\footnotesize\color{gedaempft}\addfontfeature{LetterSpace=14}\MakeUppercase{#1}}}

%% Bringt seinen Abstand selbst mit, sonst nähme er auch auf Arten ohne Pflichtsatz Raum ein. Nur für Sätze des
%% Programms, keine Eingabedaten: \programmsatz setzt den breiten Satzabstand.
\newcommand{\hinweisabsatz}[1]{\par\vspace{6mm}\noindent%
  \ungeteilt{\programmsatz{#1}}}

%% Ohne \programmsatz: Der Befreiungsgrund kommt aus den Eingabedaten (#143).
\newcommand{\steuerhinweis}[1]{\noindent\ungeteilt{#1}}

%% Mit Luft zur Tabelle, sonst klebte die erste Rechnung an der Ankündigung.
\newcommand{\absetzungskopf}[1]{\par\vspace{6mm}\noindent #1\par\vspace{1.5mm}}

%% Ein Makro samt Überschrift und Abstand, damit auf einem Dokument ohne Zahlungsaufforderung nichts stehen bleibt.
\newcommand{\bankblock}{\vspace{9mm}\noindent{\bfseries \wortbankverbindung}\\[1.5mm]%
  \wortkontoinhaber: <<KONTOINHABER>>\\
  \wortinstitut: <<INSTITUT>>\\
  IBAN: <<IBAN>>\\
  BIC: <<BIC>>}

%% Über die volle Breite: Die Zeilen tragen Nummer und zwei Daten und passen nicht in 0,55.
\newcommand{\absetzungstabelle}[1]{%
  \begin{tabularx}{\linewidth}{@{}>{\raggedright\arraybackslash}X r@{}}#1\end{tabularx}}

\newcommand{\absetzungsdetail}[1]{\newline\hspace*{4mm}{\footnotesize\color{gedaempft}#1}}

%% \newline statt \\, das in der X-Spalte die Tabellenzeile beendete. Ohne den Einzug von \absetzungsdetail;
%% der Unterschied ist Satzbild (#51).
\newcommand{\positionsdetail}[1]{\newline{\footnotesize\color{gedaempft} #1}}

%% Die Spaltenzahl steht nur hier. Die laufende Nummer setzt dieses Design nicht. Einmal gefragt, danach nur der
%% Schalter: Zeile, Spaltentyp und \spaltenzahl müssen dieselbe Antwort bekommen.
\newtoggle{rabattspalte}
\ifblank{<<RABATTSPALTE>>}{}{\toggletrue{rabattspalte}}
\newcommand{\mitrabatt}[1]{\iftoggle{rabattspalte}{#1}{}}
\iftoggle{rabattspalte}{\newcolumntype{R}{r}}{\newcolumntype{R}{}}
\newcommand{\positionszeile}[7]{#2 & #3 & #4 & #5\mitrabatt{ & #6} & #7}

%% \mbox um jeden Spaltenkopf: In der X-Spalte säße die erste Zeile sonst tiefer als in den r-/l-Spalten.
\newcommand{\positionskopf}{\mbox{\rubrik{\wortbeschreibung}} & \mbox{\rubrik{\wortmenge}} & & \mbox{\rubrik{\worteinzelpreis}}\mitrabatt{ & \mbox{\rubrik{\wortrabatt}}} & \mbox{\rubrik{\wortbetrag}} \\[1.5mm]
  \midrule}

%% Muss im ersten Argument von \multicolumn zur Zahl expandieren; \iftoggle tut das.
\newcommand{\spaltenzahl}{\iftoggle{rabattspalte}{6}{5}}

\newcommand{\summenzeile}[2]{#1 & #2}

%% Beendet die Zeile davor selbst (#147). Ein Trenner davor ginge nicht: Als \name{} vor \midrule gäbe er ein
%% „Misplaced \noalign“.
\newcommand{\abschlusszeile}[2]{\\\midrule\textbf{#1} & \textbf{#2}}

\newcommand{\nettozeile}[1]{\summenzeile{\wortnettobetrag}{#1}}
\newcommand{\steuerzeile}[3]{\summenzeile{\wortsteuer{#1}{#2}}{#3}}
\newcommand{\gesamtzeile}[1]{\abschlusszeile{\wortgesamtbetrag}{#1}}

%% ------------------------------------------------------------------ Fußzeile
%% [b], damit er nach oben wächst und nicht in den Rand. Beide Zeilen passen nachgemessen in den Satzspiegel
%% (152 und 138 mm bei 165 mm). Ohne Steuernummer: §14 Abs. 4 Nr. 2 UStG verlangt sie oder die USt-IdNr., und
%% mit beiden wurde die zweite Zeile 178 mm breit.
\newcommand{\fuss}{\parbox[b]{\textwidth}{\centering\scriptsize\color{gedaempft}
<<FIRMA>> · <<STRASSE>>, <<PLZ>> <<ORT>>\ifblank{<<REGISTERGERICHT>>}{}{ · \wortregistergericht{} <<REGISTERGERICHT>>\ifblank{<<REGISTERNUMMER>>}{}{, <<REGISTERNUMMER>>}}\\
\wortustid{} <<USTID>> · <<INSTITUT>> · BIC <<BIC>> · IBAN <<IBAN>>}}

\begin{document}
\thispagestyle{erste}

%% ------------------------------------------------------------------ Kopf
%% Das Logo hängt am Ende der ersten Zeile, denn deren Anfang hinge an der Länge der Firmierung; es nimmt keinen
%% Raum ein. \normalsize, sonst stünde Text im Logo in \small, \color{haupttext}, sonst würde eine Fläche ohne
%% eigene Farbe grau. Mit \logoskalierung liegt die Grenze bei gut 2.
%% Die Person nur ohne Geschäftsführung; bei Einzelunternehmen verlangt §14 Abs. 4 Nr. 1 UStG den Namen.
{\raggedleft\small\color{gedaempft}
  {\bfseries\color{haupttext}<<FIRMA>>}\ifblank{<<LOGO>>}{}{\settoheight{\versalhoehe}{\bfseries M}%
    \rlap{\hspace*{-\linewidth}\raisebox{\dimexpr\versalhoehe-\height\relax}[0pt][0pt]{\normalsize\color{haupttext}%
      \resizebox{!}{\logoskalierung\logohoehe}{\begin{tikzpicture}\logozeichnung\end{tikzpicture}}}}}\ifblank{<<GESCHAEFTSFUEHRUNG>>}{\kopfperson{<<INHABER>>}}{}\zeilenwechsel{}%
  <<STRASSE>> · <<PLZ>> <<ORT>>\ifblank{<<TELEFON>>}{}{\zeilenwechsel{}\worttelefon{} <<TELEFON>>}\zeilenwechsel{}<<EMAIL>>\par}

\vspace{18mm}

%% ----------------------------------------------------------- Empfängerfeld
%% Ohne Rücksendezeile, die Rechnung geht als PDF hinaus. 0.44 der Breite sind 72,6 mm, die längste Zeile braucht
%% nachgemessen 69,8 mm. Die USt-IdNr. hängt ohne Leerraum an, sonst stünde ohne sie ein Leerzeichen am Ende.
\begin{minipage}[t]{0.52\linewidth}\raggedright
  <<EMPFAENGER>>\ifblank{<<KUNDE_USTID>>}{}{\ustidtrenner{}\wortustid{} <<KUNDE_USTID>>}
\end{minipage}%
\hfill
%% Die Geschäftsführung hier und nicht im Fuß, sonst bräche die erste Fußzeile um.
\begin{minipage}[t]{0.44\linewidth}
  \raggedleft\small\color{gedaempft}
  \ifblank{<<GESCHAEFTSFUEHRUNG>>}{}{\wortgeschaeftsfuehrung\zeilenwechsel{}<<GESCHAEFTSFUEHRUNG>>\blocktrenner{}}
  \kopfangabe{\wortkundennr}{<<KUNDENNUMMER>>}
  \kopfangabe{\wortauftragsnr}{<<AUFTRAGSNUMMER>>}
  \kopfangabe{<<ZEITRAUM_BESCHRIFTUNG>>:}{\zusammen{<<LEISTUNGSZEITRAUM>>}}
  \kopfangabe{\wortvereinnahmtam}{<<VEREINNAHMT_AM>>}
  \kopfangabe{\wortreferenz}{<<KUNDENREFERENZ>>}
  \kopfangabe{\wortzahlbarbis}{<<ZAHLUNGSZIEL>>}
\end{minipage}

\vspace{14mm}

%% ---------------------------------------------------------- Titel (Blickfang)
%% Nachgemessen ist der Balken der Rechnung vor dem Generator 6,0 mm hoch. \raisebox erzwingt Ziffernhöhe und
%% Tiefe null, sonst machte die Unterlänge des „g“ in „August“ den Balken unten dicker als oben.
{\setlength{\fboxsep}{1.4mm}%
\settoheight{\balkenhoehe}{\large\bfseries 0}%
\noindent\colorbox{akzent}{%
  \raisebox{0pt}[\balkenhoehe][0pt]{%
    \makebox[\dimexpr\linewidth-2\fboxsep\relax][s]{%
      \large\bfseries\color{white}<<DOKUMENTART>> <<RECHNUNGSNUMMER>>\hfill <<RECHNUNGSDATUM>>}}}}

\vspace{9mm}

%% --------------------------------------------------------------- Einleitung
%% Aus den Eingabedaten, deshalb nicht in \hinweisabsatz (#143).
\ifblank{<<EINLEITUNG>>}{}{\noindent <<EINLEITUNG>>\par\vspace{6mm}}

%% --------------------------------------------------------- Positionen, Summen
%% Die Summen sind die letzte Zeile, angehängt mit \\*, damit sie nie ohne Position darüber stehen: Hinter der
%% Tabelle erlaubte longtable einen Umbruch, den kein \nopagebreak zurücknimmt (ADR-0020). Der Abstand als
%% \noalign, nicht als \\*[5mm], das nur die Unterlänge verlängerte; \summenluft vor der Tabelle, denn darin ist
%% \lineskip null. \LTpre und \nointerlineskip danach halten die Abstände der früheren tabularx.
\setlength{\LTpre}{\lineskip}\setlength{\LTpost}{0pt}%
\setlength{\LTleft}{0pt}\setlength{\LTright}{0pt}%
\newlength{\summenluft}\setlength{\summenluft}{\dimexpr 5mm+\lineskip\relax}%
\begin{xltabular}{\linewidth}{@{}>{\raggedright\arraybackslash}X r@{\hspace{1.5mm}}l r R r@{}}
  \positionskopf
  \endfirsthead
  \positionskopf
  \endhead
<<POSITIONEN>>\\*
  \noalign{\vskip\summenluft}
  \multicolumn{\spaltenzahl}{@{}r@{}}{\begin{minipage}[t]{0.55\linewidth}
    \begin{tabularx}{\linewidth}{@{}>{\raggedright\arraybackslash}X r@{}}
<<SUMMEN>>
    \end{tabularx}
  \end{minipage}}
\end{xltabular}
\nointerlineskip\vspace{\lineskip}

<<ARTHINWEIS>>

%% Der Abstand darunter trägt beide Fälle, mit und ohne Steuerhinweise.
\vspace{6mm}

<<STEUERHINWEISE>>

%% Umrechnung, Absetzung und Skonto auf einer Zeile, damit eine Rechnung ohne sie kein Zeichen mehr im .tex hat.
<<UMRECHNUNG>><<ABSETZUNG>><<SKONTO>>

\vspace{5mm}

%% ------------------------------------------------------------------ Schluss
%% Schlusssatz, Gruß und Name brechen nicht auseinander. Mit Strafen und nicht mit einer Box, denn Strafen
%% nehmen keinen Raum ein.
\ungeteilt{<<SCHLUSSSATZ>>}
\nopagebreak
\vspace{5mm}

\wortgruss\\*[7mm]
<<INHABER>>

%% ---------------------------------------------------------- Bankverbindung
<<BANKVERBINDUNG>>

\end{document}
